% 11/18 abstract, updated with the DPO and candidate-generator results (11/20) and the information-gain / robustness runs (11/23).
\begin{abstract}
People differ in how they want a robot to move through a crowd: some prefer it to yield early, others to keep pace;
some want it to pass behind them, others in front. Existing crowd-navigation policies optimize a fixed notion of social
compliance and cannot adapt to such individual preferences. We present a framework that identifies a user's walking-style
preference from a handful of pairwise (A/B) queries and uses it to select the robot's path at test time. A VQ-VAE trained on
human pedestrian trajectories provides a codebook of walking behaviours that generates style-diverse, human-like candidate paths.
Unsafe candidates are filtered out, and each remaining candidate is described by interpretable style features: passing side,
passing in front or behind, yielding, detour width, and human-likeness. A Bayesian Bradley--Terry model over these features is
updated after each answer, and queries are chosen by the expected utility of the best option (EUBO). On held-out situations from
Grand Central and ETH/UCY with simulated users, five queries reduce the normalized simple regret from 0.45--0.47 (human-data prior)
to 0.09--0.12, outperforming random queries in all four settings and a DPO-trained estimator on the same budget and matching an
information-gain query; the advantage over random queries persists when simulated users answer noisily or have nonlinear preferences. On shared
Grand Central scenes, codebook candidates leave 0.11 regret after five queries against 0.35 for candidates from a flow-matching
crowd-navigation generator, whose narrower candidate set bounds what any preference can achieve. In a closed-loop replay,
following the learned preference keeps the user's step regret low without raising the collision rate.
\end{abstract}
